
\documentclass[11pt,a4paper]{article}
\usepackage[utf8]{inputenc}
\usepackage[T1]{fontenc}
\usepackage[french]{babel}
\usepackage[top=3cm, bottom=2cm, left=2cm, right=2cm]{geometry}
\usepackage{stmaryrd}
\usepackage{amsmath}
\usepackage{amsfonts}
\usepackage{amssymb}
\usepackage{mathrsfs}
\usepackage{amsthm}
\usepackage{layout}
\usepackage{fancyhdr}

\newtheorem*{thm}{Théorème}
\newtheorem{ex}{Exercice}
\newtheorem*{nota}{Notation}
\newtheorem*{rem}{Remarque}
\newtheorem*{rem2}{Remarques}
\newtheorem{de2}{Définition}
\newtheorem{pro2}[de2]{Propriété}
\newtheorem{thm2}[de2]{Théorème}

\setlength{\parindent}{0cm}
\setlength{\parskip}{1ex plus 0.5ex minus 0.2ex}
\newcommand{\hsp}{\hspace{20pt}}
\newcommand{\HRule}{\rule{\linewidth}{0.5mm}}

\usepackage{comment}
\newcommand{\F}{\mathbb{F}}


\title{}

\date{}
\begin{document}


\pagestyle{fancy}

\fancyhead{}
 \fancyfoot{}

 \lhead{ 2025/2026 \\  L3 Mathématiques
}
\chead{\textbf{ Calcul formel}\\} 
 \rhead{ Université de Lorraine \\  }

\newcommand{\lb}{\llbracket}
\newcommand{\rb}{\rrbracket}
\newcommand{\KK}{\mathbb{K}}
\newcommand{\N}{\mathbb{N}}
\newcommand{\Z}{\mathbb{Z}}




\newcommand{\md}[3]{#1\ \equiv \ #2 \! \! \! \! \! \pmod {#3} }
\newcommand{\nmd}[3]{#1 \not \equiv #2 \! \! \! \! \!  \pmod {#3} }
\newcommand{\mda}[3]{#1 \equiv #2 \! \!  \pmod {#3} }
\newcommand{\nmda}[3]{#1 \not \equiv #2 \! \! \pmod {#3} }
\newcommand{\mo}[2]{#1 \! \! \! \! \! \pmod #2 }
\newcommand{\moa}[2]{#1 \! \!  \pmod {#2} }

\thispagestyle{fancy}

\begin{center}
%    \HRule \\[0.6cm]
    { \huge \bfseries
Examen
     \\ [0cm] }
    \HRule \\[0.5cm]
\end{center}






\begin{center}
\textbf{lundi 12 janvier 2026, durée : 3 heures}

\end{center}


\begin{ex}\label{monogeneite}


\begin{enumerate}

\item Soient $n\in \N_{\geq 2}$ et $G$ un sous-groupe de $(\Z/n\Z,+)$. Montrer que $G$ est monogène.

\item Déterminer les sous-groupes suivants : $\langle \overline{6},\overline{22}\rangle\subset \Z/73\Z$ et $\langle \overline{6},\overline{22}\rangle\subset \Z/74\Z$ (on précisera notamment les cardinaux de ces sous-groupes). 
\end{enumerate}

\end{ex}






\begin{ex}\label{RSA}
On s'intéresse ici à un cryptosystème utilisant l'algorithme RSA dont la clé publique vaut $(221,11)$, i.e. $n=221$ et $e=11$ : pour transmettre un message $m\in \llbracket 0,n-1\rrbracket$, on transmet le reste dans la division euclidienne de $m^e$ par $n$.
    
\begin{itemize}
\item[$1.$] Déterminer les nombres premiers $p$ et $q$ tels que $pq=221$.
\item[$2.$] Donner le chiffrement avec la clé publique du message $m=2$.
\item[$3.$] Calculer $d$ la clé privée correspondant à la clé publique $e$.
\item[$4.$]
\begin{itemize}
\item[$(a)$]
 Résoudre  l'équation suivante, d'inconnue $x\in \llbracket 0,n-1\rrbracket$:
\begin{eqnarray*}
 \left \{  \begin{array}{c}
\md{x}{3}{17},\\
\md{x}{2}{13} . \end{array}  \right. 
\end{eqnarray*}
\item[$(b)$]
En déduire le déchiffrement du message $c=7$.
\end{itemize}
\end{itemize}

\end{ex}










\begin{ex}\label{e_rac_car}

\begin{enumerate}
\item Déterminer les racines carrées de $\overline{18}$ dans $\Z/31\Z$. 

\item Déterminer les racines carrées de $\overline{18}$ dans $\Z/31^2\Z$. 


\item Déterminer les racines carrées de $\overline{45}$ dans $\Z/99\Z$. 
\end{enumerate}

\end{ex}

\begin{comment}

\begin{ex}\label{ex_divisibilite}
Soit $n\in  \Z$. 
\begin{enumerate}
\item Montrer que $6$ divise $n^3-n$.

\item Montrer que $30$ divise $n^5-n$.
\end{enumerate}
\end{ex}
\end{comment}



\begin{ex}\label{Dirichlet_faible}
Soit $X$ l'ensemble des nombres premiers congrus à $3$ modulo $4$. 
\begin{enumerate}


\item Montrer que $X$ est non-vide.

\item Montrer que l'ensemble $\{4k+1\mid k\in \N\}$ est stable par produit.

\item Soient $n\in \N^*$ et  $p_1,\ldots,p_n\in X$. Montrer que $4p_1\ldots p_n-1$ admet un diviseur dans $X$.

\item En déduire que $X$ est infini. 

\end{enumerate}
\end{ex}





\begin{ex}\label{Cipolla}[Algorithme de Cipolla (1907)]

On étudie ici l'algorithme de Cipolla, qui permet de déterminer une racine carrée d'un élément $\overline{n}$ de $\F_p=\Z/p\Z$, pour $p$ premier impair. L'algorithme fonctionne de la façon suivante. On suppose que l'on connaît  $a\in \F_p^\times$ tel que $a^2-n$ n'est pas un carré de $ \F_p$. On peut alors considérer $a+\sqrt{a^2-n}$  dans un corps contenant $\F_p$ strictement. On constate alors que $(a+\sqrt{a^2-n})^{(p+1)/2}$ est une racine carrée de $n$ dans $\F_p$. 

Les deux parties de cet exercice sont complètement indépendantes et les questions de chaque parties peuvent être traitées en admettant les précédentes.




\textbf{Première partie :}

On fixe   $p$ un nombre premier impair positif et $n\in \F_p^\times$ qui est un carré.
\begin{enumerate}
\item Soit $a\in \F_p^\times$ tel que $a^2-n$ n'est pas un carré de $\F_p$.  Soit $\mathbb{K}=\F_p[X]/(X^2-(a^2-n))$, où $X$ est une indeterminée. Montrer que $\mathbb{K}$ est un corps.

\item Déterminer un base de $\mathbb{K}$. En déduire que $\mathbb{K}$ a $p^2$ éléments.

On identifie $\F_p$ à un sous-corps de $\KK$ via l'injection $x\mapsto \overline{x}$, pour $x\in \F_p$.

\item Soit $\sqrt{a^2-n}=\overline{X}$. Que vaut $(\sqrt{a^2-n})^2$ (on justifiera rigoureusement la réponse) ?


\item Dans cette question uniquement, on choisit $p=7$, $a=\overline{1}\in \Z/7\Z$ et $n=\overline{2}\in \Z/7\Z$. On a alors $a^2-n=\overline{-1}\in \Z/7\Z$. Déterminer $(1+\sqrt{-1})^4$ dans $\KK$ (on omet les barres pour des questions de lisibilité) ? 

\item Montrer que pour tout $\alpha,\beta\in \mathbb{K}$, on a $(\alpha+\beta)^p=\alpha^p+\beta^p$.

\item Montrer que $\F_p=\{x\in \KK\mid x^p=x\}$. 

\item Montrer que $(\sqrt{a^2-n})^p=-\sqrt{a^2-n}$ (on pourra calculer $((\sqrt{a^2-n})^p)^2$).

\item Montrer que $(a+\sqrt{a^2-n})^{(p+1)/2}\in \mathbb{K}$ est une racine carrée de $n$.

\item  Montrer que $(a+\sqrt{a^2-n})^{(p+1)/2}$ est en fait un élément de $\F_p$.
\end{enumerate}

\textbf{Deuxième partie :} On fixe un nombre premier impair positif. On fixe $n\in \F_p^\times$ un carré. On veut maintenant montrer que  $\{a\in \F_p\mid a^2-n\text{ n'est pas un carré} \}$ est de cardinal $(p-1)/2$. 

\begin{enumerate}


\item Soit $a\in \F_p$. Quel calcul permet de vérifier rapidement si $a^2-n$ est un carré ? Quel est le coût en temps de ce calcul (i.e le nombre asymptotiques d'opérations élémentaires lorsque $p$ tend vers $\infty$). 

\item On pose $X=\{(a,b)\in \F_p^2\mid a^2-n=b^2\}$. Si $z\in \F_p$, on pose $X_z=\{(a,b)\in X\mid a+b=z\}$. Déterminer $X_z$ puis montrer que $|X|=p-1$. 

\textit{Indication : on pourra factoriser $a^2-b^2$.}


\item Soit $E=X\setminus (\F_p\times \{0\})$. Montrer que $|E|=p-3$ ?


\item Soit $\pi:\F_p^2\rightarrow \F_p$ définie par $\pi((x,y))=x$, pour $x\in \F_p$. Montrer que pour tout $a\in \pi(E)$, on a $|\pi^{-1}(\{a\})|=2$. 

\item Montrer que $|\pi^{-1}(a)|=1$ pour tout $a\in  \pi(X\setminus E)$.

\item Montrer que   $|\{a\in \F_p\mid a^2-n\text{ n'est pas un carré de }\F_p\}|=(p-1)/2$.


\end{enumerate}



\end{ex}

\end{document}
